% GENERATED FILE. Edit canonical lessons, then run npm run generate:books.
% canonical-source-hash: fnv1a64:0b2369cf38b1928a
% canonical-lessons: ZH-C262-wangqiu, ZH-C262-pingpangqiu, ZH-C262-duishou, ZH-C262-dabai, ZH-C262-jiayou

\chapter{Win or Lose}
\label{ch:zh-win-or-lose}
\hlchaptermodality{262}

\begin{chapteropening}
\textbf{By the end of this chapter:} \emph{I can talk about tennis and table tennis, an opponent and beating them, and cheer someone on.}

Complete the last lesson of chapter 262: I can talk about tennis and table tennis, an opponent and beating them, and cheer someone on.
\end{chapteropening}

\section[wǎngqiú]{wǎngqiú — tennis}

\label{lesson:ZH-C262-wangqiu}



\begin{quote}
Before the new one: say the Mandarin for overseas, then the Mandarin for to experience for yourself.
\end{quote}

\subsection*{You'll want to know: wǎngqiú}
\emph{wǎngqiú} — \textquotedblleft{}tennis\textquotedblright{}.

\begin{quote}
\emph{wǒ xǐhuan wǎngqiú} — I like tennis.
\end{quote}

Its characters come later; for now it is a word for the ear.

\begin{cousinweb}
\emph{wǎngqiú} is built from \zh{网} 'a net' and \emph{qiú} 'a ball', both inherited from Middle Chinese, \emph{mjangX} and \emph{gjuw} (Guangyun, in Baxter's transcription), which glosses the first 'a net' and the second 'a fine jade'; under another spelling of the second it glosses a ball of fur that is struck.
\end{cousinweb}

\begin{grammarlens}[title={What you've built}]
One more word for what I like.
\end{grammarlens}

\subsection*{Your turn}
\begin{itemize}
\raggedright
  \item \emph{Say it:} \emph{wǎngqiú}
  \item \emph{Say it:} \emph{wǎngqiú}, once more
  \item \emph{Say it:} \emph{wǒ xǐhuan wǎngqiú} — I like tennis.
  \item \emph{Recall:} say the Mandarin for overseas, then the Mandarin for to experience for yourself, then say \emph{wǎngqiú} again
\end{itemize}

\begin{tcolorbox}[breakable,colback=teal!4,colframe=teal!35!black,title={Before you move on}]
What does wǎngqiú mean? (\textquotedblleft{}Tennis\textquotedblright{}.) Say it once more.
\end{tcolorbox}
\section[pīngpāngqiú]{pīngpāngqiú — table tennis; ping-pong}

\label{lesson:ZH-C262-pingpangqiu}



\begin{quote}
Before the new one: say the Mandarin for to experience for yourself, then the Mandarin for tennis.
\end{quote}

\subsection*{You'll want to know: pīngpāngqiú}
\emph{pīngpāngqiú} — \textquotedblleft{}table tennis; ping-pong\textquotedblright{}.

\begin{quote}
\emph{tā shì pīngpāngqiú gāoshǒu} — He's a table tennis ace.
\end{quote}

Its characters come later; for now it is a word for the ear.

\begin{cousinweb}
\emph{pīngpāngqiú} is built from \emph{pīngpāng}, which imitates the sound of the ball, and \emph{qiú} 'a ball', inherited from Middle Chinese \emph{gjuw} (Guangyun, in Baxter's transcription). No dictionary we use records an older form of the first part, so its history is not settled.
\end{cousinweb}

\begin{grammarlens}[title={What you've built}]
One more word for what I like.
\end{grammarlens}

\subsection*{Your turn}
\begin{itemize}
\raggedright
  \item \emph{Say it:} \emph{pīngpāngqiú}
  \item \emph{Say it:} \emph{pīngpāngqiú}, once more
  \item \emph{Say it:} \emph{tā shì pīngpāngqiú gāoshǒu} — He's a table tennis ace.
  \item \emph{Recall:} say the Mandarin for to experience for yourself, then the Mandarin for tennis, then say \emph{pīngpāngqiú} again
\end{itemize}

\begin{tcolorbox}[breakable,colback=teal!4,colframe=teal!35!black,title={Before you move on}]
What does pīngpāngqiú mean? (\textquotedblleft{}Table tennis; ping-pong\textquotedblright{}.) Say it once more.
\end{tcolorbox}
\section[\texorpdfstring{\zh{对手} (duìshǒu)}{对手 (duìshǒu)}]{\zh{对手} (duìshǒu) — opponent; rival}

\label{lesson:ZH-C262-duishou}



\begin{quote}
Before the new one: say the Mandarin for tennis, then the Mandarin for table tennis.
\end{quote}

\subsection*{You'll want to know: \zh{对手}}
\textbf{\zh{对手}} — \emph{duìshǒu} — \textquotedblleft{}opponent; rival\textquotedblright{}.

\begin{quote}
\textbf{\zh{他是我的对手}} \emph{tā shì wǒ de duìshǒu} — He's my opponent.
\end{quote}

\begin{cousinweb}
\textbf{\zh{对手}} (\emph{duìshǒu}) is built from \zh{对} and \zh{手} 'a hand', both inherited from Middle Chinese, \emph{twojH} and \emph{syuwX} (Guangyun, in Baxter's transcription), which glosses the first 'to answer' and 'to match': the hand that matches yours.
\end{cousinweb}

\begin{grammarlens}[title={What you've built}]
One more word for what I like.
\end{grammarlens}

\subsection*{Your turn}
\begin{itemize}
\raggedright
  \item \emph{Say it:} \emph{duìshǒu}
  \item \emph{Say it:} \emph{duìshǒu}, once more
  \item \emph{Say it:} \emph{tā shì wǒ de duìshǒu} — He's my opponent.
  \item \emph{Recall:} say the Mandarin for tennis, then the Mandarin for table tennis, then say \emph{duìshǒu} again
\end{itemize}

\begin{tcolorbox}[breakable,colback=teal!4,colframe=teal!35!black,title={Before you move on}]
What does \zh{对手} mean? (\textquotedblleft{}Opponent; rival\textquotedblright{}.) Say it once more.
\end{tcolorbox}
\section[dǎbài]{dǎbài — to defeat; to beat (an opponent)}

\label{lesson:ZH-C262-dabai}



\begin{quote}
Before the new one: say the Mandarin for table tennis, then the Mandarin for an opponent.
\end{quote}

\subsection*{You'll want to know: dǎbài}
\emph{dǎbài} — \textquotedblleft{}to defeat; to beat (an opponent)\textquotedblright{}.

\begin{quote}
\emph{wǒmen dǎbàile tāmen} — We beat them.
\end{quote}

Its characters come later; for now it is a word for the ear.

\begin{cousinweb}
\emph{dǎbài} is built from \zh{打} 'to strike' and \emph{bài}, both inherited from Middle Chinese, \emph{taengX} and \emph{baejH} (Guangyun, in Baxter's transcription), which glosses the first 'to strike' and the second 'to fall to ruin by itself'. Middle Chinese \emph{taengX} would regularly give dǎng; dǎ is irregular.
\end{cousinweb}

\begin{grammarlens}[title={What you've built}]
One more word for what I like.
\end{grammarlens}

\subsection*{Your turn}
\begin{itemize}
\raggedright
  \item \emph{Say it:} \emph{dǎbài}
  \item \emph{Say it:} \emph{dǎbài}, once more
  \item \emph{Say it:} \emph{wǒmen dǎbàile tāmen} — We beat them.
  \item \emph{Recall:} say the Mandarin for table tennis, then the Mandarin for an opponent, then say \emph{dǎbài} again
\end{itemize}

\begin{tcolorbox}[breakable,colback=teal!4,colframe=teal!35!black,title={Before you move on}]
What does dǎbài mean? (\textquotedblleft{}To defeat; to beat (an opponent)\textquotedblright{}.) Say it once more.
\end{tcolorbox}
\section[jiāyóu]{jiāyóu — to cheer someone on; come on!}

\label{lesson:ZH-C262-jiayou}



\begin{quote}
Before the new one: say the Mandarin for an opponent, then the Mandarin for to defeat.
\end{quote}

\subsection*{You'll want to know: jiāyóu}
\emph{jiāyóu} — \textquotedblleft{}to cheer someone on; come on!\textquotedblright{}.

\begin{quote}
\emph{jiāyóu wǒmen néng dǎbài tāmen} — Come on, we can beat them!
\end{quote}

Its characters come later; for now it is a word for the ear.

\begin{cousinweb}
\emph{jiāyóu} is built from \emph{jiā} 'to add' and \emph{yóu} 'oil', both inherited from Middle Chinese, \emph{kae} and \emph{yuw} (Guangyun, in Baxter's transcription), which glosses the first 'to increase' and the second 'oil'. How adding oil came to mean cheering someone on is not settled.
\end{cousinweb}

\begin{grammarlens}[title={What you've built}]
One more word for what I like.
\end{grammarlens}

\subsection*{Your turn}
\begin{itemize}
\raggedright
  \item \emph{Say it:} \emph{jiāyóu}
  \item \emph{Say it:} \emph{jiāyóu}, once more
  \item \emph{Say it:} \emph{jiāyóu wǒmen néng dǎbài tāmen} — Come on, we can beat them!
  \item \emph{Recall:} say the Mandarin for an opponent, then the Mandarin for to defeat, then say \emph{jiāyóu} again
\end{itemize}

\begin{tcolorbox}[breakable,colback=teal!4,colframe=teal!35!black,title={Before you move on}]
What does jiāyóu mean? (\textquotedblleft{}To cheer someone on; come on!\textquotedblright{}.) Say it once more.
\end{tcolorbox}
